\chapter{Supplementary Code Listings}
\label{app:code}

This appendix contains supporting code listings referenced from
\cref{ch:implementation} but not essential to its main narrative.

\section{Navigation Route Definitions}

\Cref{lst:app-routes} lists the route constants used by the
\texttt{NavHost} of \cref{lst:navhost}.

\begin{lstlisting}[language=Kotlin, caption={\texttt{ui/navigation/Routes.kt}.}, label={lst:app-routes}]
object Routes {
    const val LOGIN = "login"
    const val SIGNUP = "signup"
    const val HOME = "home"
    const val ADD_REPORT = "addReport"
    const val EDIT_REPORT = "editReport/{reportId}"
    const val REPORT_DETAIL = "reportDetail/{reportId}"

    const val ARG_REPORT_ID = "reportId"

    fun editReport(reportId: Long) = "editReport/$reportId"
    fun reportDetail(reportId: Long) = "reportDetail/$reportId"
}
\end{lstlisting}

\section{Local Session Store}

\Cref{lst:app-sessionstore} shows \texttt{SessionStore}, the
\texttt{multiplatform-settings}-backed component behind OFR-2 (offline
session continuity, \cref{sec:offline-req}) and the session-resolution
logic of \cref{lst:checksession}. \texttt{Settings} is itself an
\texttt{expect}/\texttt{actual} abstraction over \texttt{SharedPreferences}
on Android and \texttt{NSUserDefaults} on iOS, so \texttt{SessionStore}
itself is common code.

\begin{lstlisting}[language=Kotlin, caption={\texttt{data/local/SessionStore.kt}.}, label={lst:app-sessionstore}]
expect fun createSettings(): Settings

class SessionStore(private val settings: Settings = createSettings()) {

    var loginTimeMillis: Long
        get() = settings.getLong(KEY_LOGIN_TIME, 0L)
        set(value) { settings[KEY_LOGIN_TIME] = value }

    var localAuthEmail: String?
        get() = settings.getStringOrNull(KEY_LOCAL_AUTH_EMAIL)
        set(value) {
            if (value != null) settings[KEY_LOCAL_AUTH_EMAIL] = value
            else settings.remove(KEY_LOCAL_AUTH_EMAIL)
        }

    fun clear() {
        settings.remove(KEY_LOGIN_TIME)
        settings.remove(KEY_LOCAL_AUTH_EMAIL)
    }

    companion object {
        const val SESSION_DURATION_MS = 60 * 60 * 1000L
        private const val KEY_LOGIN_TIME = "login_time"
        private const val KEY_LOCAL_AUTH_EMAIL = "local_auth_email"
    }
}
\end{lstlisting}

A one-hour \texttt{SESSION\_DURATION\_MS} bounds how long a user can
remain in the authenticated part of the app purely on a cached local
session before being asked to re-authenticate, balancing OFR-2's
offline-continuity goal against not caching a stale session indefinitely.

\section{Firestore Report Document Mapping}

\Cref{lst:app-remotereport} shows the exact mapping between the local
\texttt{Report} entity and the document shape uploaded to Firestore
(\cref{tab:firestore-doc}), used internally by \texttt{SyncManager}
(\cref{lst:syncmanager}).

\begin{lstlisting}[language=Kotlin, caption={Local-to-remote mapping (\texttt{sync/SyncManager.kt}).}, label={lst:app-remotereport}]
@Serializable
private data class RemoteReport(
    val localId: String,
    val title: String,
    val description: String,
    val category: String,
    val location: String,
    val photoPath: String?,
    val remotePhotoUrl: String?,
    val owner: String,
    val createdAt: Long,
    val updatedAt: Long
)

private fun Report.toRemote() = RemoteReport(
    localId = localId,
    title = title,
    description = description,
    category = category,
    location = location,
    photoPath = photoPath,
    remotePhotoUrl = remotePhotoUrl,
    owner = owner,
    createdAt = createdAt,
    updatedAt = updatedAt
)
\end{lstlisting}

Note that \texttt{id} (the local Room primary key) and \texttt{syncStatus}
are, by construction, absent from \texttt{RemoteReport} -- the type system
itself prevents these local-only bookkeeping fields from ever being
uploaded, rather than relying on a runtime filter.
